\subsubsection{Langevin diffusion and the unadjusted Langevin algorithm}
\label{subsec:mcmc-gradient-based}
\label{subsubsec:langevin-diffusion}

Many modern targets can be written in energy form,
\[
  \pi(x) \propto e^{-U(x)},
  \qquad
  x\in\R^d,
\]
where $U:\R^d\to\R$ is differentiable.
The basic random-walk Metropolis construction from \Cref{subsec:mh} treats \(\pi\) as a black box: proposals are centered at the current state and do not use any derivative information about where the density is increasing.
When gradients are available, we can do better.
The vector \(-\nabla U(x)\) points toward higher-probability regions, so it gives a local directional hint that random-walk proposals ignore.
Langevin methods use that hint while still adding noise, so the algorithm samples from the distribution instead of simply descending toward a mode.

The conceptual starting point is the overdamped Langevin diffusion
\begin{equation}
  dX_t = -\nabla U(X_t)\,dt + \sqrt{2}\,dB_t,
  \label{eq:langevin-sde}
\end{equation}
where $(B_t)_{t\ge 0}$ is standard Brownian motion in $\R^d$.
For the broader Langevin-to-SGLD viewpoint, see \citet[\S\S2--3]{welling_teh_2011_sgld}.
The drift $-\nabla U(X_t)$ points downhill in the potential and therefore toward regions where $\pi$ is larger, while the Brownian term prevents the process from collapsing onto a mode.
For small $dt$, the drift contributes an $O(dt)$ deterministic move and the noise contributes an $O(\sqrt{dt})$ random perturbation.
This combination is exactly what distinguishes sampling from optimization: the drift alone would behave like gradient descent, but the added noise keeps the process exploring the typical set rather than converging to a single maximizer.
To even formulate \eqref{eq:langevin-sde} cleanly, one usually assumes that the drift is globally Lipschitz, for example
\[
  \norm{\nabla U(x)-\nabla U(y)} \le L\norm{x-y}
  \qquad
  \text{for all }x,y\in\R^d.
\]
Under this condition the SDE is well posed: there is a unique global solution, there is no finite-time blow-up, and the law of the process depends continuously on the initial condition.
If $U\in C^2$, this is implied by a uniform Hessian bound $\norm{\Hess U(x)}\le L$.

Applying Euler--Maruyama to \eqref{eq:langevin-sde} with step size $\epsilon>0$ gives the unadjusted Langevin algorithm (ULA),
\begin{equation}
  x_{t+1}
  =
  x_t - \epsilon \nabla U(x_t) + \sqrt{2\epsilon}\,\xi_{t+1},
  \qquad
  \xi_{t+1}\sim\mathcal{N}(0,I_d),
  \label{eq:ula-update-notes}
\end{equation}
with Gaussian one-step kernel
\[
  K^{\mathrm{ULA}}_\epsilon(x,dy)
  =
  \mathcal{N}\!\left(y;\ x-\epsilon\nabla U(x),\,2\epsilon I_d\right)dy.
\]
ULA is immediately implementable and often mixes much better than a pure random walk because the drift discourages proposals deep in low-density regions.
At the same time, \eqref{eq:ula-update-notes} is only a discretization of \eqref{eq:langevin-sde}, so for fixed $\epsilon$ it usually targets an $\epsilon$-dependent approximation to $\pi$ rather than $\pi$ itself.
Reducing $\epsilon$ decreases this bias but increases autocorrelation because successive states become more similar.
The practical takeaway is therefore simple: under the standard regularity assumptions used to define the diffusion and its Fokker--Planck equation, the continuous-time Langevin diffusion preserves \(\pi\) exactly, while ULA is its cheap Euler discretization and inherits a step-size bias.
\Cref{fig:ula-energy-path} shows the drift-plus-noise geometry in one dimension.

\begin{figure}[t]
  \centering
  \begin{tikzpicture}
  \pgfplotsset{table/search path={figures/main_notes/sampling/,../../figures/main_notes/sampling/,Lecture Tex/figures/main_notes/sampling/}}
	  \begin{axis}[
	    name=energy,
	    width=0.47\linewidth,
	    height=5.8cm,
	    xmin=-3, xmax=3,
	    ymin=0, ymax=13,
	    axis lines=left,
	    xlabel={$x$},
	    ylabel={$U(x)$},
	    tick style={stat221Gray},
	    tick label style={font=\small},
	    label style={font=\small},
      legend style={
        font=\scriptsize,
        at={(axis description cs:0.5,0.98)},
        anchor=north,
        fill=white,
        fill opacity=0.90,
        draw=none,
      },
      legend columns=4,
	  ]
    \addplot[stat221Gray,very thick] table [x=x, y=U] {data_parallel_mcmc_ula_1d_energy_curve.dat};
    \addlegendentry{$U(x)$}

    \addplot[
      only marks,
      mark=*,
      mark size=1.35,
      mark options={draw=white, line width=0.25pt, fill=stat221Blue},
      opacity=0.65,
    ]
    table [x=x, y=U]
    {data_parallel_mcmc_ula_1d_energy_path.dat};
    \addlegendentry{trajectory}

    \addplot[
      only marks,
      mark=*,
      mark size=2.1,
      mark options={draw=white, line width=0.45pt, fill=stat221Purple},
    ]
    table [x=x, y=U]
    {data_parallel_mcmc_ula_1d_energy_start.dat};
    \addlegendentry{start}

    \addplot[
      only marks,
      mark=*,
      mark size=1.9,
      mark options={draw=white, line width=0.45pt, fill=stat221Teal},
    ]
    table [x=x, y=U]
    {data_parallel_mcmc_ula_1d_energy_end.dat};
    \addlegendentry{end}
  \end{axis}

	  \begin{axis}[
	    at={(energy.east)},
	    anchor=west,
	    xshift=0.06\linewidth,
	    width=0.47\linewidth,
	    height=5.8cm,
	    xmin=0, xmax=80,
	    axis lines=left,
	    xlabel={time index $t$},
	    ylabel={$x_t$},
	    tick style={stat221Gray},
	    tick label style={font=\small},
	    label style={font=\small},
	  ]
    \addplot[stat221Blue,thick,line cap=round] coordinates {(0,0.7) (1,0.6611309774)};
    \addplot[stat221Blue,thick,line cap=round] table [x=t, y=seq] {data_parallel_mcmc_ula_timeseries.dat};

    \addplot[
      only marks,
      mark=*,
      mark size=2.1,
      mark options={draw=white, line width=0.45pt, fill=stat221Purple},
    ] coordinates {(0,0.7)};
    \addplot[
      only marks,
      mark=*,
      mark size=1.9,
      mark options={draw=white, line width=0.45pt, fill=stat221Teal},
    ] coordinates {(80,1.32430078)};
  \end{axis}
\end{tikzpicture}

  \caption{A one-dimensional potential $U$ and a ULA trajectory for one fixed noise realization.
  The drift term $-\nabla U(x)$ pulls the iterate toward low-energy regions, while the Gaussian noise provides exploration.}
  \label{fig:ula-energy-path}
\end{figure}

For the continuous-time process it is useful to introduce the associated transition kernel
\[
  K_t(x,A)\coloneqq \mathbb{P}(X_t\in A\mid X_0=x),
\]
and the Markov semigroup acting on test functions,
\[
  (P_t g)(x)\coloneqq \E[g(X_t)\mid X_0=x]=\int g(y)\,K_t(x,dy).
\]
Then invariance means \(\pi P_t=\pi\) for every \(t\ge 0\), and the generator is the infinitesimal derivative
\[
  \mathcal{L}g=\lim_{t\downarrow 0}\frac{P_t g-g}{t}
\]
on smooth functions.
The kernels satisfy the continuous-time Chapman--Kolmogorov relation
\[
  K_{s+t}(x,A)=\int K_t(y,A)\,K_s(x,dy),
  \qquad
  s,t\ge 0,
\]
or equivalently \(P_{s+t}=P_sP_t\).
This semigroup viewpoint is the clean way to organize the next calculation: the generator describes the local drift-diffusion balance, while the semigroup describes how those infinitesimal moves accumulate over time.

The next calculation records the infinitesimal balance behind that exactness.
At stationarity, the inward drift and the outward diffusion cancel one another so that the target law is preserved rather than collapsed onto the mode.
For a smooth test function $g$, the generator of \eqref{eq:langevin-sde} is
\begin{equation}
  (\mathcal{L}g)(x)
  =
  -\ip{\nabla U(x)}{\nabla g(x)} + \Delta g(x),
  \label{eq:langevin-generator}
\end{equation}
and the corresponding Fokker--Planck equation for the density $f_t$ of $X_t$ is
\begin{equation}
  \partial_t f_t(x)
  =
  \nabla\cdot\bigl(f_t(x)\nabla U(x)\bigr)+\Delta f_t(x)
  =
  \nabla\cdot\bigl(f_t(x)\nabla U(x)+\nabla f_t(x)\bigr).
  \label{eq:fokker-planck-langevin}
\end{equation}
The probability flux in \eqref{eq:fokker-planck-langevin} vanishes when $f_t(x)\propto e^{-U(x)}$, because then $\nabla f_t=-f_t\nabla U$.
This is the infinitesimal reason the diffusion preserves the Gibbs density.

\begin{proposition}[Generator identity for the Langevin Gibbs law]
\label{prop:langevin-stationary}
Assume $U$ is differentiable and $e^{-U}$ is integrable on $\R^d$.
Let $\pi(x)=Z^{-1}e^{-U(x)}$.
If boundary terms vanish for the test functions of interest, then
\[
  \int (\mathcal{L}g)(x)\,\pi(x)\,dx = 0
\]
for every smooth test function $g$ with sufficient decay.
\end{proposition}
\begin{proof}
Let $g$ be smooth with sufficient decay.
Using \eqref{eq:langevin-generator},
\[
  \int (\mathcal{L}g)(x)\,\pi(x)\,dx
  =
  \frac{1}{Z}\int \left[-\ip{\nabla U(x)}{\nabla g(x)}+\Delta g(x)\right]e^{-U(x)}\,dx.
\]
Integration by parts gives
\[
  \int (\Delta g)e^{-U}\,dx
  =
  -\int \ip{\nabla g}{\nabla(e^{-U})}\,dx
  =
  \int \ip{\nabla g}{\nabla U}\,e^{-U}\,dx,
  \qquad
  \nabla(e^{-U})=-e^{-U}\nabla U,
\]
so the two terms cancel and $\int (\mathcal{L}g)\,\pi=0$.
\end{proof}
This generator identity is the infinitesimal stationarity condition.
Combined with the standard well-posedness assumptions that justify the Fokker--Planck equation and its stationary solutions, it yields invariance of \(\pi\) for the Langevin semigroup because the flux in \eqref{eq:fokker-planck-langevin} vanishes when \(f_t=\pi\).

To quantify convergence, it is convenient to work with couplings.
A \emph{coupling} of probability measures \(\mu\) and \(\nu\) on \(\R^d\) is a joint law \(\gamma\) on \(\R^d\times\R^d\) whose marginals are \(\mu\) and \(\nu\).
The 2-Wasserstein distance is then
\[
  W_2(\mu,\nu)
  \coloneqq
  \inf_{\gamma\in\mathcal{C}(\mu,\nu)}
  \left(\int \norm{x-y}^2\,\gamma(dx,dy)\right)^{1/2},
\]
where \(\mathcal{C}(\mu,\nu)\) denotes the set of all couplings.
This metric is natural for Langevin diffusions because synchronous couplings give direct pathwise control of \(\norm{X_t-Y_t}\).
Earlier, in \Cref{subsec:ergodicity}, we used total variation because it controls discrepancies of arbitrary events for general chains.
Here the dynamics have an explicit Euclidean geometry, so \(W_2\) is often the more informative metric: it is built from couplings of whole trajectories and interacts directly with convexity of the drift.

Under strong convexity the diffusion also forgets its initialization at an exponential rate.
The cleanest proof uses synchronous coupling: if two copies of the diffusion share the same Brownian motion, then the noise cancels in their difference and only the contracting drift remains.
For a probability measure \(\mu\), we write \(\mu K_t\) for the law obtained by evolving the initial distribution \(\mu\) under the Langevin diffusion for time \(t\).

\begin{proposition}[Synchronous coupling gives exponential Wasserstein contraction]
\label{prop:langevin-w2-contraction}
Assume $\nabla U$ is globally Lipschitz and that $U$ is $m$-strongly convex, meaning
\[
  \ip{\nabla U(x)-\nabla U(y)}{x-y} \ge m\norm{x-y}^2
  \qquad
  \text{for all }x,y\in\R^d.
\]
Let $(X_t)$ and $(Y_t)$ be the corresponding unique strong solutions to \eqref{eq:langevin-sde} driven by the same Brownian motion and started from $X_0=x$, $Y_0=y$.
Then for all $t\ge 0$,
\[
  \norm{X_t-Y_t}\le e^{-mt}\norm{x-y}.
\]
Consequently, for any probability measures $\mu,\nu$ on $\R^d$,
\[
  W_2(\mu K_t,\nu K_t)\le e^{-mt}W_2(\mu,\nu),
\]
where $K_t$ denotes the Langevin transition kernel at time $t$.
\end{proposition}
\begin{proof}
With the same Brownian motion, the difference $Z_t\coloneqq X_t-Y_t$ satisfies the integral equation
\[
  Z_t = Z_0 - \int_0^t \bigl(\nabla U(X_s)-\nabla U(Y_s)\bigr)\,ds.
\]
Hence $Z_t$ is absolutely continuous, and for almost every $t$,
\[
  \frac{d}{dt}\norm{Z_t}^2
  =
  -2\ip{Z_t}{\nabla U(X_t)-\nabla U(Y_t)}
  \le
  -2m\norm{Z_t}^2.
\]
Gronwall's inequality yields $\norm{Z_t}\le e^{-mt}\norm{Z_0}$.
Now let $\gamma_0$ be any coupling of $\mu$ and $\nu$, draw $(X_0,Y_0)\sim\gamma_0$, and evolve both coordinates with the same Brownian motion.
Then $(X_t,Y_t)$ is a coupling of $\mu K_t$ and $\nu K_t$, so
\[
  W_2^2(\mu K_t,\nu K_t)
  \le
  \E\norm{X_t-Y_t}^2
  \le
  e^{-2mt}\E\norm{X_0-Y_0}^2.
\]
Taking the infimum over all initial couplings $\gamma_0$ gives
\[
  W_2^2(\mu K_t,\nu K_t)\le e^{-2mt}W_2^2(\mu,\nu),
\]
and hence the claimed contraction bound.
\end{proof}

This result explains why gradient-based samplers are especially clean for strongly log-concave targets: the continuous-time dynamics contract naturally toward equilibrium.
For the Euler discretization, however, one must separate two distinct questions.
If \(K_\epsilon\) denotes the ULA kernel and \(\pi_\epsilon\) its invariant law when it exists, then
\[
  W_2(\mu K_\epsilon^n,\pi)
  \le
  W_2(\mu K_\epsilon^n,\pi_\epsilon)
  +
  W_2(\pi_\epsilon,\pi).
\]
The first term is convergence of the discrete chain toward its own equilibrium, while the second is the step-size bias introduced by discretizing the diffusion.
Under strong log-concavity and sufficiently small \(\epsilon\), ULA is itself geometrically ergodic toward \(\pi_\epsilon\); the central limitation is that at fixed \(\epsilon\) one typically has \(\pi_\epsilon\neq\pi\).
For strongly log-concave targets with \(L\)-Lipschitz gradient, this creates a genuine two-scale requirement: when \(0<\epsilon<2/L\), one can drive \(W_2(\mu K_\epsilon^n,\pi_\epsilon)\) down geometrically in \(n\), but one still has to choose \(\epsilon\) small enough that \(W_2(\pi_\epsilon,\pi)\) is itself below the target tolerance; see \citet[Thm.~1 and \S7.1]{vempala_wibisono_2019_ula_isoperimetry} for one precise local reference.

This decomposition clarifies the practical tradeoff.
Larger \(\epsilon\) makes each step more aggressive and can reduce the number of iterations needed to approach \(\pi_\epsilon\), but it also enlarges the asymptotic bias.
Smaller \(\epsilon\) reduces that bias while making neighboring states more correlated.
For ULA there is therefore no single notion of ``mixing fast'': to hit a target accuracy \(\delta\), one typically needs both
\[
  W_2(\mu K_\epsilon^n,\pi_\epsilon)\lesssim \delta/2
  \qquad\text{and}\qquad
  W_2(\pi_\epsilon,\pi)\lesssim \delta/2.
\]
The first condition is a mixing requirement for the discrete chain, while the second is an approximation requirement for the time step.

The same split appears in Monte Carlo estimation.
If we estimate \(\E_\pi[g(X)]\) by the empirical average
\[
  \hat\mu_n^{\mathrm{ULA}}
  \coloneqq
  \frac{1}{n}\sum_{k=1}^n g(X_k),
\]
then after burn-in the chain average is centered at \(\E_{\pi_\epsilon}[g(X)]\), not at \(\E_\pi[g(X)]\).
The total error therefore combines initialization bias, Monte Carlo variance from the correlated chain, and discretization bias \(\E_{\pi_\epsilon}[g]-\E_\pi[g]\).
Equivalently,
\[
  \hat\mu_n^{\mathrm{ULA}}-\E_\pi[g]
  =
  \underbrace{\hat\mu_n^{\mathrm{ULA}}-\E_{\pi_\epsilon}[g]}_{\text{initialization + correlated Monte Carlo}}
  +
  \underbrace{\E_{\pi_\epsilon}[g]-\E_\pi[g]}_{\text{discretization bias}}.
\]
For an ergodic ULA chain the first term vanishes as \(n\to\infty\), but the second persists whenever \(\epsilon\) is held fixed.
For first-order Euler discretizations that bias is typically \(O(\epsilon)\) for smooth observables, which is why it is useful to check whether estimates stabilize as the step size is reduced.
In practice one either treats ULA as an approximate sampler and checks that estimates stabilize as \(\epsilon\) is reduced, or removes the bias with a Metropolis correction.

\subsubsection{MALA, stochastic gradients, and the road to HMC}
\label{subsubsec:mala}

The Metropolis-adjusted Langevin algorithm (MALA) removes the discretization bias of ULA by using the ULA transition as a Metropolis--Hastings proposal.
With
\begin{equation}
  q_\epsilon(y\mid x)
  =
  \mathcal{N}\!\left(y;\ x-\epsilon\nabla U(x),\,2\epsilon I_d\right),
  \label{eq:mala-proposal}
\end{equation}
the acceptance probability is
\begin{equation}
  \alpha(x,y)
  =
  \min\!\left\{1,\frac{\tilde\pi(y)\,q_\epsilon(x\mid y)}{\tilde\pi(x)\,q_\epsilon(y\mid x)}\right\}
  =
  \min\!\left\{1,\exp\!\bigl(\log r(x,y)\bigr)\right\},
  \label{eq:mala-accept}
\end{equation}
where
\begin{equation}
  \log r(x,y)
  =
  -U(y)+U(x)
  -\frac{1}{4\epsilon}
  \left(
    \norm{x-y+\epsilon\nabla U(y)}^2
    -
    \norm{y-x+\epsilon\nabla U(x)}^2
  \right).
  \label{eq:mala-log-ratio}
\end{equation}
Computing \eqref{eq:mala-log-ratio} on the log scale is the stable implementation of the Metropolis correction.
By \Cref{prop:mh-detailed-balance}, the resulting kernel leaves $\pi$ invariant exactly.

The correction term in \eqref{eq:mala-log-ratio} has a simple interpretation.
The first part, $-U(y)+U(x)$, is the usual log target ratio from Metropolis--Hastings.
The quadratic part corrects for proposal asymmetry: it checks whether the forward drift from $x$ makes $y$ plausible and whether the reverse drift from $y$ makes $x$ plausible.
Small step sizes therefore make MALA behave like ULA with an acceptance rate near $1$, while larger step sizes produce more rejections because the Euler proposal no longer tracks the diffusion accurately.

\begin{algorithm}[t]
  \caption{Metropolis-adjusted Langevin algorithm (MALA)}
  \label{alg:mala}
  \begin{algorithmic}[1]
    \Require Potential $U(x)$ and gradient $\nabla U(x)$, step size $\epsilon>0$, initial state $x^{(0)}$, steps $T$
    \For{$t=0,1,\dots,T-1$}
      \State Draw $\xi\sim\mathcal{N}(0,I_d)$ and set $y \gets x^{(t)}-\epsilon\nabla U(x^{(t)})+\sqrt{2\epsilon}\,\xi$.
      \State Compute $\log r\bigl(x^{(t)},y\bigr)$ from \eqref{eq:mala-log-ratio} and draw $u\sim\mathrm{Unif}(0,1)$.
      \If{$\log u \le \min\{0,\log r\bigl(x^{(t)},y\bigr)\}$}
        \State $x^{(t+1)} \gets y$
      \Else
        \State $x^{(t+1)} \gets x^{(t)}$
      \EndIf
    \EndFor
    \State \Return $x^{(0)},x^{(1)},\dots,x^{(T)}$
  \end{algorithmic}
\end{algorithm}

\begin{figure}[t]
  \centering
  \begin{tikzpicture}[>=Stealth]
  % Level sets of a correlated Gaussian (same parameterization as fig_gibbs_2d_geometry).
  \def\a{0.810} % rho=0.95
  \def\b{0.586} % rho=0.95

  \begin{axis}[
    name=rw,
    width=0.46\linewidth,
    height=5.8cm,
    axis equal image,
    xmin=-2.6, xmax=2.6,
    ymin=-2.1, ymax=2.1,
    axis lines=middle,
    axis line style={stat221Gray!45,line width=0.6pt},
    title={random-walk proposal},
    title style={font=\scriptsize, text=stat221Gray, at={(axis description cs:0.5,1.04)}, anchor=south},
    xlabel={},
    ylabel={},
    xtick=\empty,
    ytick=\empty,
    clip=false,
  ]
    \addplot[stat221Blue!55!black, line width=0.8pt, opacity=0.35, domain=0:360, samples=240]
      ({2.0*(\a*cos(x)+\b*sin(x))},{2.0*(\b*cos(x)+\a*sin(x))});
    \addplot[stat221Blue!55!black, line width=0.8pt, opacity=0.22, domain=0:360, samples=240]
      ({1.2*(\a*cos(x)+\b*sin(x))},{1.2*(\b*cos(x)+\a*sin(x))});

    \coordinate (xrw) at (axis cs:1.85,1.35);

    \path[fill=stat221GrayLight, fill opacity=0.70] (xrw) circle[radius=0.62];
    \draw[stat221Gray!70, line width=0.9pt, dashed] (xrw) circle[radius=0.62];

    \addplot[
      only marks,
      mark=*,
      mark size=2.8pt,
      mark options={draw=white, line width=0.35pt, fill=stat221Purple},
    ] coordinates {(1.85,1.35)};
    \node[font=\scriptsize, anchor=south west, fill=white, inner sep=0.9pt] at (xrw) {$x$};

    \node[font=\scriptsize, anchor=west, fill=white, inner sep=0.9pt] at (axis cs:2.42,0.12) {$x_1$};
    \node[font=\scriptsize, anchor=south, fill=white, inner sep=0.9pt] at (axis cs:0.12,2.00) {$x_2$};
  \end{axis}

  \begin{axis}[
    at={(rw.east)},
    anchor=west,
    xshift=0.08\linewidth,
    name=lan,
    width=0.46\linewidth,
    height=5.8cm,
    axis equal image,
    xmin=-2.6, xmax=2.6,
    ymin=-2.1, ymax=2.1,
    axis lines=middle,
    axis line style={stat221Gray!45,line width=0.6pt},
    title={Langevin proposal},
    title style={font=\scriptsize, text=stat221Teal, at={(axis description cs:0.5,1.04)}, anchor=south},
    xlabel={},
    ylabel={},
    xtick=\empty,
    ytick=\empty,
    clip=false,
  ]
    \addplot[stat221Blue!55!black, line width=0.8pt, opacity=0.35, domain=0:360, samples=240]
      ({2.0*(\a*cos(x)+\b*sin(x))},{2.0*(\b*cos(x)+\a*sin(x))});
    \addplot[stat221Blue!55!black, line width=0.8pt, opacity=0.22, domain=0:360, samples=240]
      ({1.2*(\a*cos(x)+\b*sin(x))},{1.2*(\b*cos(x)+\a*sin(x))});

    \coordinate (xlan) at (axis cs:1.85,1.35);
    \coordinate (mlan) at (axis cs:1.25,0.85);

    \path[fill=stat221TealLight, fill opacity=0.70] (mlan) circle[radius=0.62];
    \draw[stat221Teal, line width=1.0pt] (mlan) circle[radius=0.62];

    \draw[->, stat221Red, line width=1.0pt] (xlan) -- (mlan)
      ;

    \addplot[
      only marks,
      mark=*,
      mark size=2.8pt,
      mark options={draw=white, line width=0.35pt, fill=stat221Purple},
    ] coordinates {(1.85,1.35)};
    \addplot[
      only marks,
      mark=*,
      mark size=2.8pt,
      mark options={draw=white, line width=0.35pt, fill=stat221Teal},
    ] coordinates {(1.25,0.85)};

    \node[
      font=\scriptsize,
      anchor=south west,
      fill=white,
      inner sep=0.9pt,
    ] at (xlan) {$x$};
    \node[
      font=\scriptsize,
      anchor=north west,
      fill=white,
      inner sep=0.9pt,
      text=stat221Teal,
    ] at (mlan) {$x-\epsilon\nabla U(x)$};

    \node[font=\scriptsize, anchor=west, fill=white, inner sep=0.9pt] at (axis cs:2.42,0.12) {$x_1$};
    \node[font=\scriptsize, anchor=south, fill=white, inner sep=0.9pt] at (axis cs:0.12,2.00) {$x_2$};
  \end{axis}
\end{tikzpicture}

  \caption{Schematic comparison of a random-walk Gaussian proposal (dashed) and a Langevin proposal (solid).
  The Langevin proposal shifts the mean by $-\epsilon\nabla U(x)$, biasing proposals toward higher-density regions.}
  \label{fig:mala-proposal-geometry}
\end{figure}

Neal's discussion of the Langevin method makes the connection to HMC explicit: one leapfrog step of Hamiltonian Monte Carlo, followed by a Metropolis correction, is exactly a Langevin proposal of this form; see \citet[\S5.2]{neal_2011_hmc}.
From that augmented-state viewpoint, MALA is the one-step overdamped relative of HMC.
HMC keeps the gradient information but carries momentum over many leapfrog steps, which suppresses the random-walk behavior that remains when every move is only a local Langevin correction.
NUTS then replaces the hand-tuned trajectory length by the no-U-turn stopping rule; see \citet[\S3.1]{hoffman_gelman_2014_nuts} and \Cref{subsec:hmc,subsec:nuts}.
This is the main conceptual bridge: ULA and MALA already show how gradients tilt local proposals, while HMC and NUTS show how augmentation converts that same local information into long coherent trajectories.

In large-data problems, evaluating the full gradient can dominate the cost of each step.
If
\[
  U(x)=\sum_{i=1}^N u_i(x),
\]
then stochastic gradient Langevin dynamics (SGLD) replaces the exact gradient by a mini-batch estimate; this is the large-data variant introduced by \citet[\S3]{welling_teh_2011_sgld}.
With batch size $B$, draw indices $i_1,\dots,i_B$ uniformly and define
\[
  \hat{\nabla U}(x)
  \coloneqq
  \frac{N}{B}\sum_{b=1}^B \nabla u_{i_b}(x).
\]
The update becomes
\begin{equation}
  x_{t+1}
  =
  x_t - \epsilon_t\hat{\nabla U}(x_t) + \sqrt{2\epsilon_t}\,\xi_{t+1},
  \qquad
  \xi_{t+1}\sim\mathcal{N}(0,I_d),
  \label{eq:sgld-update}
\end{equation}
with a user-chosen step-size schedule $(\epsilon_t)$.
If we write
\[
  \hat{\nabla U}(x_t)=\nabla U(x_t)+\eta_t,
  \qquad
  \E[\eta_t\mid x_t]=0,
\]
then
\begin{equation}
  x_{t+1}
  =
  x_t - \epsilon_t\nabla U(x_t) + \sqrt{2\epsilon_t}\,\xi_{t+1} - \epsilon_t\eta_t.
  \label{eq:sgld-two-noises}
\end{equation}
Equation \eqref{eq:sgld-two-noises} makes the approximation explicit: besides the calibrated Langevin noise, SGLD contains extra noise from subsampling.
For very small step sizes that additional term becomes negligible, but at practical step sizes it can substantially change the stationary behavior.
SGLD is therefore best understood as an approximate large-scale sampler whose per-iteration cost is cheap precisely because it gives up the exact invariance enjoyed by MALA.

Two operating regimes are worth separating.
If asymptotic exactness matters, then one uses a decreasing schedule with \(\sum_t \epsilon_t=\infty\) and \(\sum_t \epsilon_t^2<\infty\), so the Euler bias and the mini-batch perturbation are both annealed away over time.
This is the regime analyzed in the original SGLD paper \citep[\S3]{welling_teh_2011_sgld}.
In large-scale practice, however, fixed or slowly decaying step sizes are more common because they mix faster per unit gradient cost; the price is that the chain then settles near a biased limiting distribution, just as ULA does; see \citet[\S4.1]{welling_teh_2011_sgld} for the corresponding practical discussion of the transition into the Langevin phase.
This is why practitioners often use larger step sizes during early exploration or warm-up and smaller ones when collecting approximate posterior draws.

The batch size must be chosen jointly with the step size.
The stochastic-gradient perturbation in \eqref{eq:sgld-two-noises} enters at scale \(\epsilon_t\eta_t\), whereas the calibrated Langevin perturbation enters at scale \(\sqrt{\epsilon_t}\,\xi_{t+1}\).
Under sampling with replacement, the covariance of the gradient-estimation error typically scales like \(N^2\Sigma(x_t)/B\) for a suitable per-datum covariance matrix \(\Sigma(x_t)\).
Consequently the extra SGLD perturbation has covariance on the order of \(\epsilon_t^2 N^2\Sigma(x_t)/B\), whereas the injected Langevin noise has covariance \(2\epsilon_t I_d\).
Thus very noisy mini-batches can dominate the intended diffusion unless \(\epsilon_t\) is small enough or \(B\) is large enough.
A practical rule is to choose \(B\) so that the mini-batch term is a correction to the Gaussian noise rather than the main source of randomness once the chain is near stationarity; if one increases \(\epsilon_t\), one usually has to increase \(B\) as well.
This is why SGLD often uses rougher settings only for early exploration, then reduces the step size and often enlarges the batch size before collecting approximate posterior draws.
In practice it is useful to calibrate this near a high-density point rather than in the tails.
A short pilot run can estimate the scale of the mini-batch gradient error, and the diagnostic question is simple: is the extra perturbation \(\epsilon_t^2\operatorname{Var}(\eta_t\mid x_t)\) comfortably smaller than the injected Langevin variance \(2\epsilon_t I_d\)?
If not, one should shrink \(\epsilon_t\), increase \(B\), or both.
Batch sizes in the same rough range used in stochastic optimization are a common starting point, but stability checks across several \((B,\epsilon)\) pairs are more informative than trusting one nominal default; this is exactly the kind of schedule-versus-noise calibration emphasized by \citet[\S4]{welling_teh_2011_sgld}.
Comparing posterior histograms or summary estimates across those settings is often the quickest way to tell whether the stochastic-gradient noise is dominating the intended diffusion.

At a high level, the tradeoffs are now visible.
ULA is the cheapest gradient-based method but carries step-size bias, MALA removes that bias with an accept/reject step, SGLD trades exactness for mini-batch scalability, and HMC or NUTS spend full gradients to convert the same local information into much longer coherent trajectories.
Equivalently, ULA and fixed-step SGLD buy cheaper iterations by tolerating an \(\epsilon\)-dependent limiting law, whereas MALA, HMC, and NUTS spend extra work per iteration to keep the exact target invariant.

All of these gradient-based samplers inherit the transformed-density bookkeeping from \Cref{subsubsec:transformed-densities}.
If the computation runs on unconstrained coordinates $z$, then both the potential $U(z)$ and its gradient must include the Jacobian contribution from the change of variables.
Otherwise the sampler follows the right local geometry for the wrong target distribution.
